\section{Related Work}
\label{sec:related_work}

\subsection{Tiny Object Detection \& Scale-Aware Representations}
Tiny object detection in unconstrained aerial imagery, maritime surveillance, and satellite remote sensing has emerged as a distinct, critical research domain due to extreme scale variations, low signal-to-noise ratios, and severe background clutter~\cite{yu2020tinyperson,ying2025safit,lin2017fpn}. Standard object detection frameworks (e.g., Faster R-CNN~\cite{ren2015faster}, Cascade R-CNN, and RetinaNet) designed for generic benchmarks like MS COCO~\cite{lin2014coco} operate predominantly on medium-to-large instances where object regions encompass thousands of pixels. Under repeated spatial downsampling (e.g., standard stride-32 backbones), features belonging to microscopic targets ($<8\times 8$ pixels) are rapidly annihilated in deeper convolutional layers.

To mitigate scale discrepancy, several specialized paradigms have been explored. Feature Pyramid Networks (FPN)~\cite{lin2017fpn} and its extensions construct top-down semantic pathways to combine high-resolution spatial details with deep semantic features. Image tiling and patch-based cropping strategies (e.g., SAHI, SND) split high-resolution images into overlapping sub-tiles to artificially enlarge relative object scales. In addition, Yu et al.~\cite{yu2020tinyperson} introduced the TinyPerson benchmark and Scale Match, aligning feature scale distributions between source and target datasets. Similarly, Ying et al.~\cite{ying2025safit} developed visible-thermal paired benchmarks for cross-spectral tiny target reasoning. However, existing methods often require multi-stage cropping or ultra-dense anchor grids, drastically inflating inference latency and memory footprints. In contrast, our proposed framework introduces zero inference overhead by operating directly within standard multi-scale architectures.

\subsection{Distance Metrics \& Optimal Transport for Box Regression}
Bounding box regression and label assignment historically rely on discrete Intersection-over-Union (IoU) metrics. To alleviate the vanishing gradient problem of IoU on non-overlapping boxes, numerous bounding box regression losses have been proposed, including Generalized IoU (GIoU), Distance-IoU (DIoU), Complete IoU (CIoU), and Efficient IoU (EIoU). While these metrics improve convergence on standard benchmarks, they remain fundamentally dependent on intersecting bounding box areas, thus suffering from discrete step collapse when tiny objects shift by even 1--2 pixels.

To establish continuous similarity for non-overlapping bounding boxes, recent studies have explored optimal transport and Wasserstein metrics. Normalized Wasserstein Distance (NWD)~\cite{xu2022nwd} models bounding boxes as 2D Gaussian distributions and computes their 2-Wasserstein distance. Receptive Field Label Assignment (RFLA)~\cite{xu2022rfla} substitutes discrete IoU with Gaussian receptive field matching. Subsequent works including SIMD~\cite{shi2024simd}, GCD~\cite{guan2025gcd}, SWIN-TOD~\cite{wang2024swintod}, MMPW-Net~\cite{su2024mmpw}, and IGWD~\cite{hu2026igwd} introduce various geometric interpolations. 

Nonetheless, modeling bounding boxes as 2D Gaussian distributions incurs severe theoretical and empirical drawbacks. The Gaussian density approximation is inherently isotropic and spatial-smoothing, which blurs crisp object boundaries. As a direct consequence, while Gaussian metrics improve loose recall at low IoU thresholds ($\text{AP}_{50}$), they experience catastrophic performance degradation at stringent localization thresholds ($\text{AP}_{75}$). In this work, we propose the Anisotropic Log-Wasserstein (ALW) metric, which explicitly decouples center translation from aspect-ratio mismatch while normalizing scale along independent spatial axes, maintaining sharp boundary localization.

\subsection{Gradient Conflict Resolution \& Curriculum Distillation}
In multi-task and multi-scale learning architectures, conflicting gradient directions across different objectives can lead to destructive interference and negative transfer. Gradient projection techniques, such as PCGrad (Projecting Conflicting Gradients) and CAGrad, project conflicting task gradients onto orthogonal planes to prevent dominant tasks from corrupting secondary objectives. Concurrently, Curriculum-Based Learning (CBL)~\cite{chen2024dila} organizes training data or loss functions from easy to hard scales to stabilize early optimization. In this work, we adapt orthogonal gradient projection to the micro-instance proposal regime via Proposal Micro-Rescue (PC-MR) and stabilize dynamic multi-scale curriculum updates via FPN Cosine Feature Distillation (PC-MOC).
